\chapter{A toy example}

\begin{definition}[Property $A$]\label{def:A}\lean{Toy.A}
  For every $n$, $A(n)$ states that $n + 1 \le n + 2$.
\end{definition}

\begin{definition}[Property $B$]\label{def:B}\lean{Toy.B}
  For every $n$, $B(n)$ states that $n + 2 \le n + 3$.
\end{definition}

\begin{lemma}\label{lem:lemma1}\lean{Toy.lemma1}\leanok
  For every $n$, $A(n)$ holds.
  \uses{def:A}
\end{lemma}

\begin{proof}
  Trivial by arithmetic.
\end{proof}

\begin{lemma}\label{lem:lemma2}\lean{Toy.lemma2}\leanok
  For every $n$, $A(n+1)$ holds.
  \uses{def:A}
\end{lemma}

\begin{proof}
  Also trivial.
\end{proof}

\begin{lemma}\label{lem:lemma3}\lean{Toy.lemma3}\leanok
  For every $n$, $A(n) \wedge B(n+1)$ holds.
  \uses{def:B}
\end{lemma}

\begin{proof}
  Just use lemma2.
\end{proof}

\begin{lemma}\label{lem:main}\lean{Toy.main}\leanok
  For every $n$, $A(n) \wedge B(n+1)$ holds.
  \uses{lem:lemma1,lem:lemma3}
\end{lemma}

\begin{proof}
  Easy from lemma1 and lemma3.
\end{proof}
